\documentclass[12pt,a4paper]{article}

% Preamble generated by build_preamble for exercises

% Source: sources/preamble/exercises_packages.tex

\usepackage[T1]{fontenc}
\usepackage[utf8]{inputenc}
\usepackage[british]{babel}
\usepackage{csquotes}

\usepackage{geometry}
\geometry{
  a4paper,
  left=25mm,
  right=25mm,
  top=25mm,
  bottom=30mm
}

\usepackage{mathtools}
\usepackage{amsthm}

\usepackage{libertine}
\usepackage[libertine]{newtxmath}
\usepackage[mathscr]{euscript}

\usepackage{enumitem}
\usepackage{tikz-cd}
\usepackage{comment}

\usepackage{xcolor}
\definecolor{CustomBlue}{RGB}{51,51,153}

\usepackage{hyperref}
\hypersetup{
  colorlinks=true,
  linkcolor=CustomBlue,
  citecolor=CustomBlue,
  urlcolor=CustomBlue
}
\urlstyle{same}

\usepackage[noabbrev]{cleveref}

% Source: sources/preamble/exercises_macros.tex

% Exercise sheet metadata
\newcounter{sheet}
\setcounter{sheet}{1}

\newcommand{\sheetcourse}{Course}
\newcommand{\sheetlecturer}{Lecturer}
\newcommand{\sheetduedate}{Due date}

\newcommand{\course}[1]{\renewcommand{\sheetcourse}{#1}}
\newcommand{\sheetnumber}[1]{\setcounter{sheet}{#1}}
\newcommand{\lecturer}[1]{\renewcommand{\sheetlecturer}{#1}}
\newcommand{\duedate}[1]{\renewcommand{\sheetduedate}{#1}}

\newcommand{\makesheetheader}{%
    \hypersetup{
        pdftitle={\sheetcourse\space --- Sheet \arabic{sheet}},
        pdfauthor={\sheetlecturer}
    }%
    \noindent
    \begin{tabular*}{\textwidth}{@{}l@{\extracolsep{\fill}}c r@{}}
        \textbf{Lecturer:} & \textbf{\sheetcourse} & \textbf{Due date:} \\
        \sheetlecturer & Sheet \arabic{sheet} & \sheetduedate
    \end{tabular*}%
    \par\vspace{0.6em}%
    \noindent\hrule
    \vspace{1.2em}%
}

% Exercise environments
\theoremstyle{definition}
\newtheorem{exercise}{Exercise}
\renewcommand{\theexercise}{\arabic{sheet}.\arabic{exercise}}

\crefname{exercise}{exercise}{exercises}
\Crefname{exercise}{Exercise}{Exercises}

% Exercise parts
\setlist[enumerate,1]{
    label=(\alph*),
    ref=(\alph*)
}

\crefname{enumi}{part}{parts}
\Crefname{enumi}{Part}{Parts}

% Solutions
\newif\ifshowsolutions
\showsolutionsfalse

\NewDocumentEnvironment{solution}{+b}
{%
    \ifshowsolutions
        \begingroup
            \renewcommand{\qedsymbol}{$\blacksquare$}%
            \begin{proof}[Solution]
                #1
            \end{proof}
        \endgroup
    \fi
}
{}

% Source: sources/preamble/todo.tex

% Warning symbol
\usepackage{utfsym}

% To-do notes
\newif\ifshowtodos
\showtodosfalse

\newcommand{\todo}[1]{%
    \ifshowtodos
        \noindent{\color{red}\usym{26A0} #1}%
    \fi
}

% Source: config/user_example.sh

\course{Course}
\sheetnumber{1}
\lecturer{Example Author}
\duedate{Due date}

% Source: sources/examples/exercises.tex

\begin{document}

\course{Algebra I}
\sheetnumber{3}
\duedate{15 October 2026}

\showtodostrue

\makesheetheader

\begin{exercise}
    \label{exercise:basic}
    Let \(V\) be a finite-dimensional vector space.
    \begin{enumerate}
        \item Define a canonical map \(V \to V^{\vee\vee}\).
            \label{item:dual-dual-map}
        \item Prove that this map is an isomorphism.
            \label{item:dual-dual-isomorphism}
    \end{enumerate}
\end{exercise}

\showsolutionstrue
\begin{solution}
    For \(v \in V\), define
    \[
        \Phi(v) \colon V^\vee \to k,
        \qquad
        \Phi(v)(\lambda)=\lambda(v).
    \]
    This gives a linear map \(\Phi \colon V \to V^{\vee\vee}\), which is the map
    asked for in \Cref{item:dual-dual-map}. After choosing a basis of \(V\), the
    map \(\Phi\) is identified with the usual identification of a finite-dimensional
    vector space with its double dual. Hence \(\Phi\) is an isomorphism, proving
    \Cref{item:dual-dual-isomorphism}.
\end{solution}
\showsolutionsfalse

\begin{exercise}
    \label{exercise:exact-sequence}
    Let
    \[
        0 \to A^{\vee} \to B^{\vee} \to C^{\vee} \to 0
    \]
    be a short exact sequence of finite-dimensional vector spaces. Use
    \Cref{item:dual-dual-isomorphism} of \Cref{exercise:basic} to describe the
    dual exact sequence.
\end{exercise}

\begin{solution}
    Dualizing gives an exact sequence
    \[
        0 \to C^{\vee\vee} \to B^{\vee\vee} \to A^{\vee\vee} \to 0.
    \]
    The finite-dimensionality assumption allows us to identify each vector space
    with its double dual, using \Cref{item:dual-dual-isomorphism} of
    \Cref{exercise:basic}.
    So we obtain a short exact sequence
    \[
        0 \to C \to B \to A \to 0.
    \]
\end{solution}

\end{document}
